% Generated draft: requires mathematical and source review.
\documentclass[11pt]{article}
\usepackage[margin=25mm]{geometry}
\usepackage{fontspec}
\setmainfont{Latin Modern Roman}
\usepackage{amsmath,amssymb,mathtools,unicode-math}
\setmathfont{Latin Modern Math}
\usepackage{enumitem,needspace,xurl,hyperref}
\setlength{\parindent}{0pt}
\setlength{\parskip}{5pt}
\setlength{\emergencystretch}{3em}
\newcommand{\N}{\mathbb{N}}
\newcommand{\Z}{\mathbb{Z}}
\newcommand{\Q}{\mathbb{Q}}
\newcommand{\R}{\mathbb{R}}
\newcommand{\C}{\mathbb{C}}
\newcommand{\F}{\mathbb{F}}
\newcommand{\A}{\mathbb{A}}
\newcommand{\PP}{\mathbb{P}}
\newcommand{\E}{\mathbb{E}}
\newcommand{\eps}{\varepsilon}
\newcommand{\dd}{\,\mathrm{d}}
\begin{document}
\section*{UnsolvedMath ID: 6000004}
\textbf{Source problem code:} catalogue entry
\subsection{Definitions and mathematical statement}
A sequence of random variables $Z_n$ satisfies a large-deviation principle with speed $a_n\to\infty$ and lower-semicontinuous rate $I$ if the lower bound on open sets and upper bound on closed sets have the form
\[
 -\inf_{x\in O}I(x)\leq\liminf_n\frac1{a_n}\log\Pr(Z_n\in O),
 \qquad
 \limsup_n\frac1{a_n}\log\Pr(Z_n\in C)\leq-\inf_{x\in C}I(x).
\]
Relative entropy of probabilities is
\[
 D(P\|Q)=\int\log\left(\frac{dP}{dQ}\right)\,dP
 \quad(P\ll Q),\qquad D(P\|Q)=+\infty\text{ otherwise}.
\]
The source concerns large-deviation settings with relative-entropy rate functions; not every large-deviation rate is asserted to have this form.

On an information manifold with metric $g$, two affine connections $\nabla,\nabla^*$ are dual when
\[
 Xg(Y,Z)=g(\nabla_XY,Z)+g(Y,\nabla_X^*Z)
\]
for vector fields $X,Y,Z$. Torsion, flatness, choice of statistical model, and any $\alpha$-connection parameter require additional hypotheses.

\textbf{Research theme (Shun-ichi Amari).} Investigate how relative-entropy large-deviation principles relate to dual-connection structures in information geometry. In particular, seek precise links between their probabilistic rate functions/divergences and the geometric metric/connection data. The excerpt does not prescribe a class of probability models, a rate-function construction, a specific dual pair, or a theorem to establish, so the theme is not a uniquely quantified universal assertion.

The supplied summary says the relative-entropy/canonical-divergence connection is well developed, refers to dual or $\alpha$-connection structures, and leaves a fully general exhaustive theory open/in flux. It gives no precise theorem or regularity hypotheses. This does not assert that arbitrary dual-connection structures are realized by entropy-rate LDPs, nor that one relative-entropy divergence realizes every parameter choice. The record's open label is a research-theme status rather than a declaration that the basic connection is unproved.
\subsection{Short English statement}
Study how large-deviation principles with relative-entropy rate functions encode dual affine connections in information geometry. The basic link is reported as established, but the source asks a broad research theme without specifying models, connection classes, or a target theorem.
\subsection{Sources}
\begin{itemize}
\item Local UnsolvedMath registry, permanent ID 6000004, source code catalogue entry, ``Large Deviations and Dual Connections''; proposer: Shun-ichi Amari, proposed year 1998. Status open, classification OPEN-TRIAGE, updated 2026-08-06.
\item Furuhata, Matsuzoe, and Urakawa, \emph{Open Problems in Affine Differential Geometry and Related Topics} (1998), item 1(d), scanned PDF pages 1--2, \url{https://www.jstage.jst.go.jp/article/iis/4/2/4_2_125/_pdf/-char/en}. Reported manual OCR from page image and access: 2026-07-29. The article presents the item as open; current status NEEDS\_REVIEW. The record says the J-STAGE article page marks it CC BY 4.0.
\item The summary describes substantial literature and canonical-divergence links without identifying specific publications or exhaustive realization results.
\end{itemize}
\end{document}
